\section{Parent-complete terminal-faithful pair lifts}
\label{sec:tpc64-parent-faithful-lifts}

A pair relation can be useful only if its rectangular lift represents the
same terminal field without deleting the parent against which activation
is measured.  This section gives the exact existence test, the unique
least relation, and the resulting parent inflation.

\subsection{Block rectangles and two active closures}

Fix \(k\in K_X\).  For a relation
\[
 I\subseteq\mathcal M_X\times\mathcal R_{X,k},
\]
define its block rectangle inside the declared ambient incidence by
\begin{equation}
 \operatorname{Rect}_k(I)
 :=
 \{u\in\Omega_{X,k}:(m(u),r(u))\in I\}.
 \label{eq:tpc64-block-rectangle}
\end{equation}
Its pre-terminal and terminal-active closures are
\begin{align}
 \operatorname{PRect}_k(I)
 &:=
 \{u\in\operatorname{Rect}_k(I):c_u\ne0\},
 \label{eq:tpc64-block-pre-rectangle}\\
 \operatorname{TRect}_{h,k}(I)
 &:=
 \{u\in\operatorname{Rect}_k(I):t_h(u)c_u\ne0\}.
 \label{eq:tpc64-block-terminal-rectangle}
\end{align}
Both are monotone in \(I\).

\begin{definition}[Parent-complete and terminal-atomic faithful]
\label{def:tpc64-parent-complete-terminal-faithful}
A pair relation \(I\) on block \(k\) is
\begin{enumerate}[label=\textup{(\roman*)},leftmargin=3em]
 \item \emph{parent-complete} if
 \begin{equation}
  \mathcal U_{X,k}^{\mathrm{pre},+}
  \subseteq\operatorname{PRect}_k(I);
  \label{eq:tpc64-parent-complete}
 \end{equation}
 \item \emph{terminal-atomic faithful} if
 \begin{equation}
  \operatorname{TRect}_{h,k}(I)
  =\mathcal U_{X,h,k}^{+}.
  \label{eq:tpc64-terminal-atomic-faithful}
 \end{equation}
\end{enumerate}
\end{definition}

Parent completeness retains every nonzero literal pre-terminal coordinate.
Terminal-atomic faithfulness demands equality before any coherent grouping;
terminal-active spurious crosses are therefore forbidden even when a later
grouping might cancel them.

Put
\begin{equation}
 I_{\mathrm{pre},k}
 :=
 \operatorname{pr}_{m,r}
 \mathcal U_{X,k}^{\mathrm{pre},+}.
 \label{eq:tpc64-minimal-parent-relation}
\end{equation}
By construction,
\begin{equation}
 I\text{ is parent-complete}
 \quad\Longleftrightarrow\quad
 I_{\mathrm{pre},k}\subseteq I.
 \label{eq:tpc64-parent-complete-relation-inclusion}
\end{equation}

\subsection{The exact existence and minimality theorem}

\begin{theorem}[Parent-complete faithful-lift criterion]
\label{thm:tpc64-parent-complete-faithful-lift}
There exists a pair relation on block \(k\) which is simultaneously
parent-complete and terminal-atomic faithful if and only if
\begin{equation}
 \boxed{
 \operatorname{TRect}_{h,k}(I_{\mathrm{pre},k})
 =\mathcal U_{X,h,k}^{+}.}
 \label{eq:tpc64-parent-complete-lift-iff}
\end{equation}
When \eqref{eq:tpc64-parent-complete-lift-iff} holds,
\(I_{\mathrm{pre},k}\) is the unique least such relation.  More precisely,
\(I\) is parent-complete and terminal-atomic faithful if and only if
\begin{equation}
 I_{\mathrm{pre},k}\subseteq I
 \quad\text{and}\quad
 \operatorname{TRect}_{h,k}
     (I\setminus I_{\mathrm{pre},k})=\varnothing.
 \label{eq:tpc64-all-parent-complete-faithful-relations}
\end{equation}
\end{theorem}

\begin{proof}
Every point of \(\mathcal U_{X,h,k}^{+}\) belongs to
\(\mathcal U_{X,k}^{\mathrm{pre},+}\), so its pair lies in
\(I_{\mathrm{pre},k}\).  Hence
\begin{equation}
 \mathcal U_{X,h,k}^{+}
 \subseteq
 \operatorname{TRect}_{h,k}(I_{\mathrm{pre},k}).
 \label{eq:tpc64-terminal-contained-in-parent-closure}
\end{equation}
If \(I\) is parent-complete, then
\(I_{\mathrm{pre},k}\subseteq I\) by
\eqref{eq:tpc64-parent-complete-relation-inclusion}.  If it is also
terminal-atomic faithful, monotonicity gives
\[
 \mathcal U_{X,h,k}^{+}
 \subseteq
 \operatorname{TRect}_{h,k}(I_{\mathrm{pre},k})
 \subseteq
 \operatorname{TRect}_{h,k}(I)
 =\mathcal U_{X,h,k}^{+},
\]
which proves necessity of
\eqref{eq:tpc64-parent-complete-lift-iff}.  Conversely, that equality
makes \(I_{\mathrm{pre},k}\) both parent-complete and faithful.
Every parent-complete relation contains it, proving unique leastness.

Assume the criterion.  Adding pairs outside \(I_{\mathrm{pre},k}\)
preserves terminal equality exactly when those pairs create no
terminal-active point, which is the second condition in
\eqref{eq:tpc64-all-parent-complete-faithful-relations}.  The first
condition is exactly parent completeness.
\end{proof}

If \eqref{eq:tpc64-parent-complete-lift-iff} fails, every pair relation
which retains all literal parent atoms creates a terminal-active spurious
cross.  No enlargement of the relation can remove that cross.  Hence no
parent-complete, terminal-atomic faithful pair-rectangular lift exists.  The
canonical repair used below then falls back to the literal active
hypergraph.  This statement excludes pair rectangles only; it makes no
minimality claim among arbitrary nonlinear, quotient, or grouped
representations.

\subsection{Exact parent inflation}

For a parent-complete relation \(I\), define its populated rectangular
row weight by
\begin{equation}
 \widetilde P_{k,m}(I)
 :=
 \sum_{\substack{u\in\operatorname{PRect}_k(I)\\m(u)=m}}
 |c_u|^2
 \label{eq:tpc64-rectangular-parent-weight}
\end{equation}
and its added row mass by
\begin{equation}
 \Delta_{k,m}(I)
 :=
 \sum_{\substack{
 u\in\operatorname{PRect}_k(I)
       \setminus\mathcal U_{X,k}^{\mathrm{pre},+}\\
 m(u)=m}}
 |c_u|^2.
 \label{eq:tpc64-block-added-parent-mass}
\end{equation}
Because parent completeness gives inclusion of the literal active
incidence, the union is disjoint and
\begin{equation}
 \boxed{
 \widetilde P_{k,m}(I)
 =P_{k,m}+\Delta_{k,m}(I).}
 \label{eq:tpc64-block-parent-inflation-identity}
\end{equation}

Likewise define
\begin{equation}
 \widetilde T_{k,m}(I)
 :=
 \sum_{\substack{
 u\in\operatorname{TRect}_{h,k}(I)\\m(u)=m}}
 |t_h(u)c_u|^2.
 \label{eq:tpc64-rectangular-terminal-weight}
\end{equation}
If \(I\) is terminal-atomic faithful, then
\begin{equation}
 \boxed{\widetilde T_{k,m}(I)=T_{k,m}}
 \label{eq:tpc64-block-terminal-preservation}
\end{equation}
for every row.  Thus a terminal-faithful rectangle can change the
denominator while leaving the terminal numerator exactly unchanged.

\begin{theorem}[Minimal parent inflation and its zero test]
\label{thm:tpc64-minimal-parent-inflation}
Assume \eqref{eq:tpc64-parent-complete-lift-iff}.  Then, for every
parent-complete terminal-atomic faithful relation \(I\),
\begin{equation}
 \Delta_{k,m}(I_{\mathrm{pre},k})
 \le\Delta_{k,m}(I)
 \qquad(m\in\mathcal M_X).
 \label{eq:tpc64-rowwise-minimal-inflation}
\end{equation}
The canonical pair lift has zero parent inflation on every row if and only
if the literal pre-terminal active incidence is rectangular over its pair
projection:
\begin{equation}
 \boxed{
 \Delta_{k,m}(I_{\mathrm{pre},k})=0
 \text{ for every }m
 \quad\Longleftrightarrow\quad
 \operatorname{PRect}_k(I_{\mathrm{pre},k})
 =\mathcal U_{X,k}^{\mathrm{pre},+}.}
 \label{eq:tpc64-zero-inflation-iff}
\end{equation}
If the terminal criterion holds but
\eqref{eq:tpc64-zero-inflation-iff} fails, then some row has strictly
positive added parent mass, and every added nonzero coordinate is
terminal-null.
\end{theorem}

\begin{proof}
Every admissible \(I\) contains \(I_{\mathrm{pre},k}\).  Monotonicity of
\(\operatorname{PRect}_k\), followed by termwise nonnegativity in
\eqref{eq:tpc64-block-added-parent-mass}, proves
\eqref{eq:tpc64-rowwise-minimal-inflation}.  The inclusion
\[
 \mathcal U_{X,k}^{\mathrm{pre},+}
 \subseteq\operatorname{PRect}_k(I_{\mathrm{pre},k})
\]
always holds.  Therefore all rowwise added sums vanish exactly when the
two finite active sets are equal, proving
\eqref{eq:tpc64-zero-inflation-iff}.  If they are unequal, a spurious
point has \(c_u\ne0\) and contributes positively to its row.  Terminal
faithfulness forces \(t_h(u)c_u=0\) at every such point.
\end{proof}

For \(P_{k,m}>0\), define the block inflation
\begin{equation}
 \iota_{k,m}(I)
 :=\frac{\widetilde P_{k,m}(I)}{P_{k,m}}
 =1+\frac{\Delta_{k,m}(I)}{P_{k,m}}.
 \label{eq:tpc64-block-inflation-factor}
\end{equation}
If \(P_{k,m}=\widetilde P_{k,m}(I)=0\), put
\(\iota_{k,m}(I)=1\).  If
\(P_{k,m}=0<\widetilde P_{k,m}(I)\), put
\(\iota_{k,m}(I)=+\infty\).  The last case cannot occur for the canonical
relation \(I_{\mathrm{pre},k}\), but it can occur after arbitrary extra
rows are inserted.

The identities
\eqref{eq:tpc64-block-parent-inflation-identity} and
\eqref{eq:tpc64-block-terminal-preservation} show exactly which
denominator a later activation theorem must use.  They do not by
themselves give a positive activation constant.

\begin{remark}[Terminal-only least relations do not preserve the parent]
\label{rem:tpc64-terminal-only-relation}
The smaller relation
\(\operatorname{pr}_{m,r}\mathcal U_{X,h,k}^{+}\) may represent the
terminal atomic field while omitting pairs which carry nonzero
pre-terminal mass.  A favorable ratio obtained after that deletion is not
a lower bound relative to \(P_{k,m}\).  Parent completeness is precisely
the condition preventing this substitution.
\end{remark}

The exact size of the resulting activation dilution, new-only row
failures, and sharp finite countermodels are separate from the existence
theorem.  They must be evaluated from the populated repaired parent, not
from the number of pair or block labels.
